We now apply the Fredholm alternative theorem (Theorem~\ref{thm:fredholm-alternative-compact}) to the operator $G$ to get the following.

\begin{Theorem}
\label{thm:elliptic-fredholm-alternative}
Let $L(\bm{x},D)$ be a uniformly elliptic differential operator of order $2k$ satisfying the hypotheses of Theorem~\ref{thm:weak-dirichlet-positive-shift-existence}. Then for every $\mu\in\mathbb{C}$ the Fredholm alternative holds; i.e., either\begin{enumerate}\item for every $f\in L^2(\Omega)$ there exists a unique weak solution $u\in H_0^k(\Omega)$ of the Dirichlet problem for the equation
\begin{equation}\label{eq:elliptic-fredholm-inhomogeneous-equation}
L(\bm{x},D)u-\mu u=f,
\end{equation}
i.e.,
\begin{equation}\label{eq:elliptic-fredholm-inhomogeneous-weak-equation}
B[v,u]-\mu(v,u)_{L^2(\Omega)}=(v,f)_{L^2(\Omega)}
\end{equation}
for all $v\in H_0^k(\Omega)$, or\item there exists at most a finite linearly independent collection of functions $u_i\in H_0^k(\Omega)$, $i=1,\ldots,N$, such that
\begin{equation}\label{eq:elliptic-fredholm-homogeneous-weak-equations}
B[v,u_i]-\mu(v,u_i)=0,
\end{equation}
for all $v\in H_0^k(\Omega)$.\end{enumerate}Furthermore, the set of values at which the second alternative holds forms an infinite discrete set with no finite accumulation point.
\end{Theorem}

\begin{Proof}
We first write the equation
\begin{equation}\label{eq:elliptic-fredholm-differential-eigenproblem}
Lu=\mu u
\end{equation}
as
\begin{equation}\label{eq:elliptic-fredholm-shifted-eigenproblem}
(L+\widetilde\lambda)u=(\widetilde\lambda+\mu)u.
\end{equation}
Then by a formal calculation in which we act on both sides of \eqref{eq:elliptic-fredholm-shifted-eigenproblem} with $G/\widetilde\lambda=(L+\widetilde\lambda)^{-1}$ we see that \eqref{eq:elliptic-fredholm-differential-eigenproblem} holds if and only if $u$ solves
\begin{equation}\label{eq:elliptic-fredholm-solution-operator-eigenrelation}
u=(L+\widetilde\lambda)^{-1}(L+\widetilde\lambda)u=\dfrac{\widetilde\lambda+\mu}{\widetilde\lambda}Gu=\dfrac{1}{\sigma}Gu.
\end{equation}
Thus, we see that $u\in L^2(\Omega)$ is an eigenfunction of $G$ corresponding to the eigenvalue $\sigma$ if and only if $u$ is an eigenfunction of $L$ corresponding to the eigenvalue $\mu$ where
\begin{equation}\label{eq:elliptic-fredholm-sigma-from-mu}
\sigma=\dfrac{\widetilde\lambda}{\widetilde\lambda+\mu},
\end{equation}

\begin{equation}\label{eq:elliptic-fredholm-mu-from-sigma}
\mu=-\widetilde\lambda+\dfrac{\widetilde\lambda}{\sigma}.
\end{equation}
By the Fredholm alternative theorem, the nonzero eigenvalues of $G$ are of finite multiplicity and thus the eigenvalues of $L$ are as well. Also, the eigenvalues of $G$ form a discrete set whose only possible accumulation point is zero, and since we have arranged it so that $0$ is not an eigenvalue of $G$, $G$ must have an infinite collection of eigenvalues. Thus, there must be an infinite collection of eigenvalues of $L$ with no finite accumulation point.

When $\mu\ne-\widetilde\lambda$ is not an eigenvalue of $L$, we note that $u\in H_0^k(\Omega)$ is a solution of \eqref{eq:elliptic-fredholm-inhomogeneous-equation} if and only if $u$ is a solution of
\begin{equation}\label{eq:elliptic-fredholm-equivalent-compact-equation}
G(u)-\dfrac{\widetilde\lambda}{\mu+\widetilde\lambda}u=-\dfrac{1}{\mu+\widetilde\lambda}G(f).
\end{equation}
We leave it to the reader to supply the rigor necessary to shore up this formal argument. The only delicate points involve showing that functions $u$ that are solutions of equations involving $G$ (and are thus naturally thought of as being only in $L^2(\Omega)$) must actually be functions in $H_0^k(\Omega)$ imbedded into $L^2(\Omega)$ (and can thus work as weak solutions of equations involving $L$).
\end{Proof}
